\clearpage
\VerdierPhysicalPageAnchor{0051}%
\VerdierSemanticAnchor{ch1-1-5}
\subsection{\texorpdfstring{$\mathbb Z^n$}{Z n}-范畴}

\VerdierSemanticAnchor{ch1-1-5-1}
\noindent\textbf{1.5.1.}\quad
设 $(\mathcal C,\mathbb Z^n,T,c)$ 为一个 $\mathbb Z^n$-范畴，
$(\mathcal C',\mathbb Z^n,T',c')$ 为一个加法 $\mathbb Z^n$-范畴，并设
\[
  (F,p)\colon(\mathcal C,\mathbb Z^n,T,c)
  \longrightarrow(\mathcal C',\mathbb Z^n,T',c')
\]
为一个由正规化余循环
$\alpha\in Z^2(\mathbb Z^n,\mathbb Z/2\mathbb Z)$ 扭曲的
$\mathbb Z^n$-函子。以 $e_1,\ldots,e_n$ 表示 $\mathbb Z^n$ 的典范基。
于是，对每个 $i\in[1,n]$，都有函子同构

\VerdierSemanticAnchor{ch1-1-5-1-eq1}
\begin{equation}
  p(e_i)\colon FT(e_i)\longrightarrow T'(e_i)F
  \tag{1.5.1.1}
\end{equation}
并且对每一对 $i\ne j$，下图为
$\alpha(e_i,e_j)\alpha^{-1}(e_j,e_i)$-交换的：

\VerdierSemanticAnchor{ch1-1-5-1-eq2}
\begin{center}
  \resizebox{0.98\textwidth}{!}{%
    $\begin{tikzcd}[ampersand replacement=\&,column sep=4.2em,row sep=3.3em]
      FT(e_i+e_j)
        \arrow[r,"{F*c(e_i,e_j)}"]
        \arrow[d,"{F*c(e_j,e_i)}"]
        \&
      FT(e_j)T(e_i)
        \arrow[r,"{p(e_j)*T(e_i)}"]
        \&
      T'(e_j)FT(e_i)
        \arrow[d,"{T'(e_j)*p(e_i)}"'] \\
      FT(e_i)T(e_j)
        \arrow[d,"{p(e_i)*T(e_j)}"]
        \& {} \&
      T'(e_j)T'(e_i)F
        \arrow[d,"{c'(e_i,e_j)^{-1}*F}"'] \\
      T'(e_i)FT(e_j)
        \arrow[r,"{T'(e_i)*p(e_j)}"']
        \&
      T'(e_i)T'(e_j)F
        \arrow[r,"{c'(e_j,e_i)^{-1}*F}"']
        \&
      T'(e_i+e_j)F
    \end{tikzcd}$%
  }\;.
  \hfill\textup{(1.5.1.2)}
\end{center}

\VerdierSemanticAnchor{ch1-1-5-2}
\noindent\textbf{命题 1.5.2.}\quad
\emph{设 $F\colon\mathcal C\to\mathcal C'$ 为一个函子，并且对每个
$i\in[1,n]$，给定函子同构
\[
  p_i\colon FT(e_i)\longrightarrow T'(e_i)F,
\]
使得对每一对 $i\ne j$，在图 (1.5.1.2) 中把 $p(e_i)$ 换成 $p_i$
所得的图为
$\alpha(e_i,e_j)\alpha^{-1}(e_j,e_i)$-交换的。则存在唯一一个由余循环
$\alpha$ 扭曲的 $\mathbb Z^n$-函子
\[
  (F,p)\colon(\mathcal C,\mathbb Z^n,T,c)
  \longrightarrow(\mathcal C',\mathbb Z^n,T',c')
\]
满足
\[
  p(e_i)=p_i,
  \qquad i\in[1,n].
\]}

\VerdierSemanticAnchor{ch1-1-5-2-proof-intro}
唯一性的证明可由考察图 (1.4.3.11) 得到；这一点留给读者验证。
我们只给出存在性证明的一些说明。证明对整数 $n$ 作归纳。
